%
% 31-ring.tex -- Funktionen auf einem Ringgebiet
%
% (c) 2026 Prof Dr Andreas Müller
%
\subsection{Funktionen auf einem Ringgebiet}
\input{chapters/040-analytisch/fig/fig-ring.tex}%
Wir betrachten eine holomorphe Funktion $f\colon U\to\mathbb{C}$.
Die offene Menge $U$ sei so gestaltet, dass der \emph{Ring}
\[
R(z_0;r_1,r_2)
=
\{
z\in\mathbb{C}
\mid
r_1 \le |z-z_0| \le r_2
\}
\]
in $U$ enthalten ist.
\index{Ringgebiet}%
Abbildung~\ref{buch:analytisch:laurent:fig:ring} zeigt diese Situation.
Da die Funktion $f$ im Inneren des kleinen Kreises mit Radius $r_1$ um
den Punkt $z_0$ nicht überall definiert sein muss, ist der Integralsatz
von Cauchy nicht auf den grossen Kreis anwendbar.
Der folgende Satz zeigt dass es aber trotzdem möglich ist, den Wert von
$f$ in einem beliebigen Punkt des Ringgebietes mit einem Wegintegral
zu berechnen.

\begin{satz}
Seien $\gamma_2$ und $\gamma_1$ die im Gegenuhrzeigersinn durchlaufenen
Randkreise des Ringgebietes $R(z_0;r_1,r_2)$ mit Radien $r_2$ und $r_1$.
Dann gilt
\begin{equation}
f(a)
=
\frac{1}{2\pi i}
\oint_{\gamma_1}
\frac{f(z)}{z-a}
\,dz
-
\frac{1}{2\pi i}
\oint_{\gamma_2}
\frac{f(z)}{z-a}
\,dz
\label{buch:analytisch:laurent:eqn:zweiintegrale}
\end{equation}
für jeden Punkt $a\in R(z_0;r_1,r_2)$.
\end{satz}

\begin{proof}
Der in Abbildung~\ref{buch:analytisch:laurent:fig:ring} dargestellte
Weg, der vom Punkt $A$ ausgehend zunächst dem äusseren Randkreis
folgt bis wieder zurück zum Punkt $A$, dann radial zum Punkt $a$, in 
einem Halbkreis vom Radius $\varepsilon$ um den Punkt $a$ bis zum
Punkt $B$, von dort im Uhrzeigersinn um den Inneren Randkreis zurück
zum Punkt $B$ und erneut radial zum Punkt $A$, mit einem kleinen, 
kreisförmigen Umweg vom Radius $\varepsilon$ um den Punkt $a$.
Die Funktion ist im Inneren des Weges überall holomorph,
nach dem Integralsatz von Cauchy gilt daher
\[
\oint_{\gamma_1-\gamma_2-\gamma_{\varepsilon}} \frac{f(z)}{z-a}\,dz 
=
0.
\]

Die radialen Teile werden zweimal in entgegengesetzter Richtung durchlaufen
und heben sich daher weg.
Einen Beitrag zum Integral leisten daher nur die Kreise.
Der kleine Kreis um den Punkt $a$ wird mit $t\mapsto \varepsilon e^{2\pi it}$
parametrisiert.
So wird das Integral
\begin{align*}
0
&=
\oint_{\gamma_1} \frac{f(z)}{z-a}\,dz
-
\oint_{\gamma_2} \frac{f(z)}{z-a}\,dz
-
\int_0^1
\frac{f(a+\varepsilon e^{2\pi it})}{a+\varepsilon e^{2\pi it}-a}
\cdot
\frac{d}{dt}
\varepsilon
e^{2\pi it}
\,dt
\\
\oint_{\gamma_1} \frac{f(z)}{z-a}\,dz
-
\oint_{\gamma_2} \frac{f(z)}{z-a}\,dz
&=
\int_0^1
\frac{f(a+\varepsilon e^{2\pi it})}{\varepsilon e^{2\pi it}}
\cdot
2\pi i\varepsilon e^{2\pi it}
\,dt
\\
&=
2\pi i
\int_0^1
f(a+\varepsilon e^{2\pi it})
\,dt.
\end{align*}
Die rechte Seite geht für $\varepsilon\to 0$ gegen $0$, so dass
\[
f(a)
=
\frac{1}{2\pi i}
\oint_{\gamma_1}
\frac{f(z)}{z-a}
\,dz
-
\frac{1}{2\pi i}
\oint_{\gamma_2}
\frac{f(z)}{z-a}
\,dz
\]
folgt.
\end{proof}
